\documentclass{beamer}
\usepackage{amsfonts,amsmath,oldgerm}
\usepackage{booktabs}
\usepackage{array}
\usepackage{graphicx}
\usepackage{tabularx}
\usepackage{xeCJK}
\usetheme{sintef}

\usefonttheme[onlymath]{serif}
\titlebackground*{assets/background}

\newcommand{\metric}[3]{\texttt{#1 / #2 / #3}}
\newcommand{\metriccell}[3]{\begin{tabular}{@{}c@{}}#1\\#2\\#3\end{tabular}}
\newcommand{\imgdir}{../images}
\newcommand{\smalltab}{\scriptsize}
\newcommand{\tight}{\vspace{-0.35em}}
\setlength{\tabcolsep}{3pt}
\renewcommand{\arraystretch}{1.08}

\title{实验课四：多模态大语言模型}
\subtitle{InternVL2-2B 监督微调、数据构造与显式 Grounding}
\author{CS60004-LAB4-MLLM}
\date{2024年5月}

\begin{document}
\maketitle

\begin{frame}[shrink=4]{实验目标与任务边界}
\small
\begin{columns}[T,onlytextwidth]
\begin{column}{0.53\textwidth}
\begin{block}{核心任务}
\begin{itemize}
\item 围绕 GQA VQA：输入 \texttt{image + question}，输出自然语言短答案。
\item 完整经历模型结构理解、SFT 训练、合成数据构造、评测与案例分析。
\item 在线 Leaderboard 最多 5 次提交，测试集只含 \texttt{image/question/bboxes}。
\end{itemize}
\end{block}
\end{column}
\begin{column}{0.43\textwidth}
\begin{block}{模型设置}
\begin{itemize}
\item 基座：\textbf{InternVL2-2B Base}
\item 教师：\textbf{Qwen3-VL-8B-Instruct}
\item 约束：单张 32GB GPU 完成训练。
\end{itemize}
\end{block}
\end{column}
\end{columns}
\vspace{0.3em}
\begin{center}
\smalltab
\begin{tabular}{p{0.18\textwidth}p{0.39\textwidth}p{0.34\textwidth}}
\toprule
任务 & 可使用数据 & BBox 规则 \\
\midrule
Task 1 & GQA image/question/answer & 不允许使用 bbox \\
Task 2 & LVIS/COCO 非 QA 数据，可与 GQA 合并 & bbox 仅作构造中间信息，最终输入不能显式保留 \\
Task 3 & GQA、Task2 数据、GQA bbox & 允许显式用于训练或推理 \\
\bottomrule
\end{tabular}
\end{center}
\end{frame}

\section{Task 1：MLLM 微调策略探索}

\begin{frame}[shrink=4]{从 Baseline 出发：Prompt 也是性能瓶颈}
\small
\begin{itemize}
\item Task 1 只允许使用 GQA 的 \texttt{image + question + answer}，不使用 bbox 与 LVIS/COCO。
\item 早期实验暴露关键现象：精确匹配评测不仅考察视觉理解，也高度依赖答案格式。
\item 在已经跑到 Task 3 后发现 prompt 对 InternVL2-2B 影响巨大，于是从 Task 1 开始整体重跑。
\end{itemize}
\vspace{0.2em}
\smalltab
\begin{tabular}{p{0.42\textwidth}ccc}
\toprule
设置 & Total & BBox & No-BBox \\
\midrule
InternVL2-2B, 早期 prompt & 0.247 & 0.2487 & 0.2222 \\
Qwen3-VL-8B, 早期 prompt & 0.685 & 0.6777 & 0.7937 \\
InternVL2-2B, 短答案 prompt & \textbf{0.685} & 0.6756 & 0.8254 \\
Qwen3-VL-8B, 短答案 prompt & 0.679 & 0.6692 & 0.8254 \\
\bottomrule
\end{tabular}
\vspace{0.4em}
\begin{block}{Prompt 修正}
\scriptsize
旧：\texttt{Question: \{question\}\textbackslash nAnswer with the shortest correct answer only.}\\
新：\texttt{Question: \{question\}\textbackslash nAnswer with the shortest correct answer only. The answer must be a simple short answer, usually one word, a short phrase, yes/no, or a number.}
\end{block}
\end{frame}

\begin{frame}{格式错误案例：正确但无法精确匹配}
\small
\begin{columns}[T,onlytextwidth]
\begin{column}{0.45\textwidth}
\includegraphics[width=\textwidth,height=0.62\textheight,keepaspectratio]{\imgdir/task1_format_error_2386907.jpg}
\end{column}
\begin{column}{0.52\textwidth}
\begin{block}{样例}
\begin{itemize}
\item 图像：\texttt{images/gqa/2386907.jpg}
\item 问题：\textit{Is the grass brown or green?}
\item 标准答案：\texttt{green}
\item 模型输出：\texttt{The grass is green.}
\end{itemize}
\end{block}
\begin{block}{观察}
语义正确但答案不够短，精确匹配记为错误；因此后续统一使用规范化短答案 prompt。
\end{block}
\end{column}
\end{columns}
\end{frame}

\begin{frame}[shrink=4]{模块冻结实验设计}
\small
\begin{itemize}
\item InternVL2-2B 拆为视觉编码器、MLP Connector 和语言解码器。
\item 固定数据、prompt 与训练流程，仅改变可训练模块，比较效果、效率和显存。
\end{itemize}
\vspace{0.2em}
\smalltab
\begin{tabular}{p{0.10\textwidth}p{0.31\textwidth}rrp{0.26\textwidth}}
\toprule
配置 & 可训练模块 & 参数量 & 占比 & 实验含义 \\
\midrule
A & Connector & 12.60M & 0.5710\% & 只学习跨模态对齐与格式适配 \\
B & Connector + Language & 1.90B & 86.2173\% & 标准 instruction tuning 风格 \\
C & Vision + Connector & 316.61M & 14.3537\% & 视觉表征适配，不更新语言模型 \\
D & Vision + Connector + Language & 2.21B & 100.0000\% & 全参数微调，作为性能上限候选 \\
\bottomrule
\end{tabular}
\vspace{0.5em}
\begin{center}
\includegraphics[width=0.78\textwidth,height=0.26\textheight,keepaspectratio]{\imgdir/task1_wandb.png}
\end{center}
\end{frame}

\begin{frame}{主结果：全参数最优，但收益不线性}
\small
\begin{itemize}
\item 模型：InternVL2-2B；数据：GQA VQA；不使用 bbox；规范化短答案 prompt。
\item 统一设置：\texttt{batch size = 16}，\texttt{learning rate = 5e-6}，\texttt{512 training samples}。
\item 评估：验证集 \texttt{n=1000}，bbox subset 仅作事后分析。
\end{itemize}
\smalltab
\begin{tabular}{p{0.20\textwidth}cccccc}
\toprule
配置 & Total & BBox & No-BBox & Trainable & Mem & Time \\
\midrule
A Connector & 0.684 & 0.6756 & 0.8095 & 0.5710\% & 18.38G & 141.9s \\
B Conn.+Lang. & 0.701 & 0.6926 & 0.8254 & 86.2173\% & 19.39G & 176.4s \\
C Vision+Conn. & 0.688 & 0.6798 & 0.8095 & 14.3537\% & 20.41G & 229.5s \\
D Full & \textbf{0.707} & \textbf{0.6980} & \textbf{0.8413} & 100\% & 21.66G & 282.0s \\
\bottomrule
\end{tabular}
\vspace{0.3em}
\begin{block}{结论}
A 只训练 0.5710\% 参数仍达到 0.684，性价比高；B 接近 D，说明语言侧 instruction tuning 对短答案 VQA 贡献稳定；C 未超过 B，当前样本规模下单独更新视觉编码器并不更有效。
\end{block}
\end{frame}

\begin{frame}[shrink=4]{Infra 与探索性实验}
\small
\begin{columns}[T,onlytextwidth]
\begin{column}{0.44\textwidth}
\begin{block}{Infra 修改}
\begin{itemize}
\item D\_full 开启 \texttt{gradient\_checkpointing}，避免 OOM。
\item 使用梯度累积灵活设置 batch size。
\item 图片预处理改为 CPU/GPU 并行，512 samples 从 297s 降至 282s。
\item 开启 Flash Attention 后降至 186s，节省约 60\% 训练时间。
\end{itemize}
\end{block}
\end{column}
\begin{column}{0.54\textwidth}
\smalltab
\begin{tabular}{p{0.33\textwidth}rrp{0.23\textwidth}}
\toprule
设置 & Total & Time & 解释 \\
\midrule
bs128 lr4e-5 & 0.588 & 588s & 学习率偏大 \\
bs64 lr2e-5 & 0.642 & 585s & 有提升但不足 \\
bs32 lr1e-5 & 0.691 & 590s & 进入较优区间 \\
bs16 lr5e-6 & 0.695 & 593s & 小 batch 略优 \\
bs16 lr5e-6, 512 & \textbf{0.699} & 297s & 样本减半不降反升 \\
bs8 lr1e-6, 256 & 0.664 & 151s & 样本过少下降 \\
\bottomrule
\end{tabular}
\end{column}
\end{columns}
\vspace{0.3em}
\begin{block}{经验}
充分的 infra 与 prompt 探索收益很大；早期没有优先检查 prompt，导致后续实验需要整体重跑。
\end{block}
\end{frame}

\section{Task 2：从视觉标注数据构造 Instruction Tuning 数据}

\begin{frame}[shrink=5]{数据构造：从非 QA 标注到 Instruction}
\small
\begin{itemize}
\item Task 2 使用 COCO caption 与 LVIS 物体类别、bbox 标注，将非 QA 数据转化为 instruction-response 样本。
\item LVIS bbox 可作为教师模型中间提示，但最终训练与验证输入中不能显式保留 bbox 坐标或 bbox 列表。
\end{itemize}
\smalltab
\begin{tabular}{p{0.15\textwidth}p{0.28\textwidth}p{0.23\textwidth}p{0.24\textwidth}}
\toprule
方案 & 教师输入 & 输出类型 & 清洗重点 \\
\midrule
A text teacher & caption + objects + bbox 文本 & caption / 识别 / 计数 / 摘要 / 空间 & 去坐标泄漏、过滤非法 task type、去重 \\
B vision teacher & image + caption + objects + bbox 文本 & 同 A，但教师可看图纠偏 & 同 A，额外关注答案复杂度和视觉噪声 \\
\bottomrule
\end{tabular}
\vspace{0.4em}
\begin{block}{生成约束}
\begin{itemize}
\item 每张图最多生成 8 条样本，覆盖 \texttt{image\_captioning}、\texttt{object\_recognition}、\texttt{object\_counting}、\texttt{dense\_visual\_summary}、\texttt{spatial\_reasoning}。
\item 输出 JSON array，每条含 \texttt{task\_type/question/answer}；question 与 answer 使用英文。
\item 识别、计数和空间理解样本尽量做成短答案 VQA，适配精确匹配评测。
\end{itemize}
\end{block}
\end{frame}

\begin{frame}[fragile,shrink=5]{A/B 教师 Prompt 差异}
\small
\begin{columns}[T,onlytextwidth]
\begin{column}{0.49\textwidth}
\begin{block}{方案 A：纯文字描述驱动}
\tiny
\begin{verbatim}
Objects in the image:
{fmt_objects(row, keep_bbox=True)}
Image caption: {row.get("caption", "")}

Please generate 8 diverse multimodal
instruction-response samples.
\end{verbatim}
\end{block}
\end{column}
\begin{column}{0.49\textwidth}
\begin{block}{方案 B：图文联合驱动}
\tiny
\begin{verbatim}
Objects in the image:
{fmt_objects(row, keep_bbox=True)}
Image caption: {row.get("caption", "")}

Please look at the image and use the
annotations as hints. Generate 8 diverse
multimodal instruction-response samples.
\end{verbatim}
\end{block}
\end{column}
\end{columns}
\vspace{0.2em}
\begin{block}{共同规则}
\tiny
\begin{verbatim}
1. Output a JSON array. Each item must contain task_type, question, and answer.
2. task_type must be one of the allowed task types. Prefer diverse types.
3. For object_recognition/counting/spatial_reasoning, write a VQA question
   and make answer a short answer: one word, phrase, yes/no, or number.
4. The final question must not contain bbox coordinates or bbox lists.
5. Output JSON only. Do not add explanations.
\end{verbatim}
\end{block}
\end{frame}

\begin{frame}[shrink=4]{数据质量：A/B 差异主要在视觉细节可靠性}
\small
\begin{columns}[T,onlytextwidth]
\begin{column}{0.49\textwidth}
\includegraphics[width=\textwidth,height=0.36\textheight,keepaspectratio]{\imgdir/task2_case_118838_labeled.jpg}
\smalltab
\begin{tabular}{p{0.17\textwidth}p{0.70\textwidth}}
\toprule
来源 & 示例与观察 \\
\midrule
A & \textit{How many gravestones are visible?} $\to$ \texttt{3}，标注转写稳定。 \\
B & 生成含 \textit{grass/surroundings} 的描述，看图后细节更丰富。 \\
\bottomrule
\end{tabular}
\end{column}
\begin{column}{0.49\textwidth}
\includegraphics[width=\textwidth,height=0.36\textheight,keepaspectratio]{\imgdir/task2_case_388903_labeled.jpg}
\smalltab
\begin{tabular}{p{0.17\textwidth}p{0.70\textwidth}}
\toprule
来源 & 示例与观察 \\
\midrule
A & \textit{How many suitcases?} $\to$ \texttt{1}，答案形式更保守。 \\
B & \textit{How many apples?} $\to$ \texttt{many}，更真实但不适配精确匹配。 \\
\bottomrule
\end{tabular}
\end{column}
\end{columns}
\vspace{0.25em}
\begin{block}{生成规模}
A text teacher：256 images $\times$ 8，保留 1789 条；B vision teacher：保留 1522 条。B 的视觉证据更丰富，但问题和答案也更容易变长、变复杂。
\end{block}
\end{frame}

\begin{frame}{Task 2 训练结果：没有稳定超过 GQA Baseline}
\small
\begin{itemize}
\item 模型：InternVL2-2B，D\_full 全参数微调。
\item 设置：\texttt{bs16 lr5e-6 1024 samples}，后续全部用新 prompt 重跑。
\item 对比：Task2 only、Task1+Task2 混合、Task1 baseline；混合时控制总样本数为 1024。
\end{itemize}
\smalltab
\begin{tabular}{p{0.25\textwidth}ccccp{0.23\textwidth}}
\toprule
训练数据 & Total & BBox & No-BBox & $\Delta$ & 解释 \\
\midrule
Task1 baseline & \textbf{0.704} & \textbf{0.6926} & \textbf{0.8730} & 0.000 & GQA 原生监督最强 \\
A only & 0.691 & 0.6830 & 0.8095 & -0.013 & 合成数据不能替代 GQA \\
B only & 0.687 & 0.6788 & 0.8095 & -0.017 & 仍有目标分布偏差 \\
A+B & 0.691 & 0.6830 & 0.8095 & -0.013 & 数量增加但质量未升 \\
Task1+A & 0.697 & 0.6884 & 0.8254 & -0.007 & 混合后略降 \\
Task1+B & 0.699 & 0.6905 & 0.8254 & -0.005 & B 更接近 baseline \\
Task1+A+B & 0.695 & 0.6862 & 0.8254 & -0.009 & 多源噪声累积 \\
\bottomrule
\end{tabular}
\end{frame}

\begin{frame}{Task 2 结论：主要瓶颈是分布迁移}
\small
\begin{columns}[T,onlytextwidth]
\begin{column}{0.48\textwidth}
\begin{block}{正向结果}
\begin{itemize}
\item 非 QA 标注可批量生成 instruction-response。
\item 数据覆盖 caption、识别、计数、摘要和空间关系。
\item B 方案能引入更多真实视觉细节。
\end{itemize}
\end{block}
\end{column}
\begin{column}{0.48\textwidth}
\begin{block}{问题}
\begin{itemize}
\item 合成样本与 GQA 短答案 VQA 的问题风格不同。
\item 答案粒度、视觉 grounding 可靠性与验证集不一致。
\item 追求丰富 instruction 会引入 caption/summary 类分布偏差。
\end{itemize}
\end{block}
\end{column}
\end{columns}
\vspace{0.4em}
\begin{block}{后续尝试与动机}
我们又生成了更偏 GQA-style short-answer VQA 的数据，但仍没有有效提升；因此 Task 3 转向显式使用 bbox grounding，而不是继续增加弱监督合成 QA。
\end{block}
\end{frame}

\section{Task 3：Bounding Box 信息的利用}

\begin{frame}[shrink=4]{从隐式视觉理解到显式 Grounding}
\small
\begin{itemize}
\item Task 1/2 的共同瓶颈是定位敏感题：模型需要先找到正确对象，再回答颜色、属性、数量或空间关系。
\item Task 3 允许显式使用 GQA bbox，因此将问题转化为“给模型额外 grounding 信号”。
\item 验证集中 \textbf{937/1000} 个样本含 bbox，收益主要反映在 bbox subset。
\end{itemize}
\begin{columns}[T,onlytextwidth]
\begin{column}{0.42\textwidth}
\smalltab
\begin{tabular}{rrr}
\toprule
BBox 数量 & 样本数 & 占比 \\
\midrule
0 & 63 & 6.3\% \\
1 & 394 & 39.4\% \\
2 & 451 & 45.1\% \\
3 & 91 & 9.1\% \\
4 & 1 & 0.1\% \\
\bottomrule
\end{tabular}
\end{column}
\begin{column}{0.55\textwidth}
\begin{block}{探索方案}
\begin{itemize}
\item 文本化 bbox：类别 + 坐标拼接到问题前。
\item 图像可视化：同色框、彩色框、彩色框 + 标签。
\item bbox-aware prompt：颜色锚点、坐标和标签组合。
\item hard case 再训练：基于训练集错误样本定向补强。
\end{itemize}
\end{block}
\end{column}
\end{columns}
\end{frame}

\begin{frame}[fragile,shrink=8]{Prompt 侧方案}
\small
\begin{block}{no\_bbox}
\tiny
\begin{verbatim}
Question: Is the mattress to the right or to the left of the table the lamp is on?
Answer with the shortest correct answer only.
\end{verbatim}
\end{block}
\begin{block}{bbox\_prompt}
\tiny
\begin{verbatim}
Objects in the image:
1. lamp: normalized_bbox=[0.172, 0.397759, 0.276, 0.661064]
2. table: normalized_bbox=[0.174, 0.616246, 0.41, 0.935574]
3. mattress: normalized_bbox=[0.348, 0.492997, 0.944, 0.966387]

Question: Is the mattress to the right or to the left of the table the lamp is on?
Answer with the shortest correct answer only.
\end{verbatim}
\end{block}
\begin{block}{color\_bbox\_prompt}
\tiny
\begin{verbatim}
1. red box: lamp, normalized_bbox=[...]
2. blue box: table, normalized_bbox=[...]
3. green box: mattress, normalized_bbox=[...]
\end{verbatim}
\end{block}
\end{frame}

\begin{frame}{图像侧方案：bbox 可视化}
\small
\begin{columns}[T,onlytextwidth]
\begin{column}{0.32\textwidth}
\includegraphics[width=\textwidth,height=0.48\textheight,keepaspectratio]{\imgdir/task3_0953938_same_box.jpg}
\centering\scriptsize \texttt{same\_box}\\
同色框提示区域，但难区分多目标。
\end{column}
\begin{column}{0.32\textwidth}
\includegraphics[width=\textwidth,height=0.48\textheight,keepaspectratio]{\imgdir/task3_0953938_color_box.jpg}
\centering\scriptsize \texttt{color\_box}\\
彩色框支持颜色引用和多目标区分。
\end{column}
\begin{column}{0.32\textwidth}
\includegraphics[width=\textwidth,height=0.48\textheight,keepaspectratio]{\imgdir/task3_0953938_color_box_label.jpg}
\centering\scriptsize \texttt{color\_box\_label}\\
彩色框 + 编号/类别标签，语义最显式。
\end{column}
\end{columns}
\vspace{0.4em}
\begin{block}{设计取舍}
可视化把位置信息变成视觉提示，但可能遮挡原图或引入噪声；标签往往比框本身更关键。
\end{block}
\end{frame}

\begin{frame}[shrink=8]{Prompt × 图像二维消融：InternVL2-2B}
\small
\begin{itemize}
\item 不训练模型，直接比较不同推理输入形式，判断 bbox 信息本身是否有效。
\item 每个单元为 \texttt{Total / BBox / No-BBox}；最稳定的是 \texttt{bbox\_prompt + original image}。
\end{itemize}
\tiny
\begin{tabular}{p{0.16\textwidth}cccc}
\toprule
Prompt \textbackslash Image & original & same\_box & color\_box & color\_box\_label \\
\midrule
no\_bbox & \metriccell{0.685}{0.6756}{0.8254} & \metriccell{0.668}{0.6574}{0.8254} & \metriccell{0.670}{0.6596}{0.8254} & \metriccell{0.796}{0.7940}{0.8254} \\
bbox\_prompt & \textbf{\metriccell{0.799}{0.7972}{0.8254}} & \metriccell{0.781}{0.7780}{0.8254} & \metriccell{0.781}{0.7780}{0.8254} & \metriccell{0.790}{0.7876}{0.8254} \\
color\_prompt & \metriccell{0.798}{0.7962}{0.8254} & \metriccell{0.782}{0.7791}{0.8254} & \metriccell{0.772}{0.7684}{0.8254} & \metriccell{0.779}{0.7759}{0.8254} \\
color\_bbox & \metriccell{0.798}{0.7962}{0.8254} & \metriccell{0.758}{0.7535}{0.8254} & \metriccell{0.756}{0.7513}{0.8254} & \metriccell{0.790}{0.7876}{0.8254} \\
\bottomrule
\end{tabular}
\vspace{0.35em}
\begin{block}{阶段性结论}
\small
标签/文本上下文是核心：\texttt{color\_box\_label} 在 \texttt{no\_bbox} 下有效；已有 bbox prompt 时，额外画框通常不增益，甚至造成视觉噪声。
\end{block}
\end{frame}

\begin{frame}{教师模型参考：强模型更会用显式 Grounding}
\small
\begin{center}
\smalltab
\begin{tabular}{p{0.42\textwidth}cccp{0.25\textwidth}}
\toprule
模型与输入 & Total & BBox & No-BBox & 观察 \\
\midrule
Qwen3-VL-8B, no\_bbox & 0.679 & 0.6692 & 0.8254 & 缺少定位仍受限 \\
Qwen3-VL-8B, bbox\_prompt & 0.842 & 0.8431 & 0.8254 & bbox subset 大幅提升 \\
Qwen3-VL-8B, color\_bbox + label & \textbf{0.844} & \textbf{0.8453} & 0.8254 & 接近 prompt 上限 \\
\bottomrule
\end{tabular}
\end{center}
\vspace{0.4em}
\begin{block}{对后续训练的启发}
InternVL2-2B 与 Qwen3-VL-8B 都能从 bbox prompt 获益，但强模型利用更充分；因此训练主线选择 \texttt{bbox\_prompt}，可视化方案作为消融分析。
\end{block}
\end{frame}

\begin{frame}{bbox-aware 训练：模块冻结结果}
\small
\begin{itemize}
\item 在 bbox prompt 推理有效的基础上，训练目标变为让模型稳定适配 bbox-aware 输入格式。
\item 先比较 A/B/C/D 冻结配置，再探索样本量、学习率、可视化输入与 hard case。
\end{itemize}
\smalltab
\begin{tabular}{p{0.17\textwidth}ccccccp{0.17\textwidth}}
\toprule
配置 & Total & BBox & No-BBox & Trainable & Mem & Time & 观察 \\
\midrule
A Connector & 0.855 & 0.8581 & 0.8095 & 0.5710\% & 19.17G & 274s & 轻量对齐也能利用 bbox \\
B Conn.+Lang. & 0.869 & 0.8719 & 0.8254 & 86.2173\% & 19.36G & 376s & 语言侧适配关键 \\
C Vision+Conn. & 0.853 & 0.8538 & 0.8413 & 14.3537\% & 21.20G & 499s & 视觉侧收益较小 \\
D Full & \textbf{0.872} & \textbf{0.8751} & 0.8254 & 100\% & 21.63G & 603s & 后续强基线 \\
\bottomrule
\end{tabular}
\vspace{0.3em}
\begin{block}{观察}
D\_full 最高，但 B\_connector\_language 非常接近；相比视觉侧更新，语言侧对 bbox 文本格式的适配更重要。
\end{block}
\end{frame}

\begin{frame}{训练方式探索：文本 bbox 是核心}
\small
\scriptsize
\begin{tabular}{p{0.36\textwidth}cccp{0.29\textwidth}}
\toprule
实验 & Total & BBox & No-BBox & 解释 \\
\midrule
no\_bbox 训练, bbox 推理 & 0.832 & 0.8324 & 0.8254 & 仅推理加 bbox 已明显超过 Task1 \\
Task1/D\_full + bbox 512 & 0.867 & 0.8687 & 0.8413 & 继续训练有效 \\
same\_box & 0.681 & 0.6702 & 0.8413 & 纯视觉同色框失败 \\
color\_box & 0.686 & 0.6756 & 0.8413 & 彩色框仍不足 \\
color\_box\_label & 0.822 & 0.8218 & 0.8254 & 标签带来明显改善 \\
same\_box + bbox\_prompt & 0.860 & 0.8613 & 0.8413 & 文本 bbox 是核心 \\
color\_box + color\_prompt & 0.852 & 0.8527 & 0.8413 & 颜色 prompt 有效但低于 bbox prompt \\
color\_prompt without bbox & 0.860 & 0.8613 & 0.8413 & 标签/颜色锚点比坐标更关键 \\
4096 samples, bs64 lr5e-6 & 0.874 & 0.8762 & 0.8413 & 扩样小幅提升但成本高 \\
\bottomrule
\end{tabular}
\vspace{0.4em}
\begin{block}{结论}
训练与推理阶段的 bbox 表达一致性很重要；单独画框不稳定，文本化对象标签与 bbox prompt 最稳。
\end{block}
\end{frame}

\begin{frame}[fragile]{COT 与 Hard Case 尝试}
\small
\begin{columns}[T,onlytextwidth]
\begin{column}{0.50\textwidth}
\begin{block}{COT 结果}
\begin{itemize}
\item 参考 VoCoT，让模型先基于 bbox 生成区域描述和推理路径，再输出短答案。
\item 100 个训练样本上：直接 \texttt{bbox\_prompt} 为 \metric{0.840}{0.8404}{0.8333}。
\item COT 为 \metric{0.810}{0.8085}{0.8333}，输出复杂度增加，不适合短答案精确匹配。
\end{itemize}
\end{block}
\end{column}
\begin{column}{0.47\textwidth}
\begin{block}{Hard Case}
\begin{itemize}
\item 用 \texttt{task1/D\_full + bbox\_prompt} 跑 8000 条训练数据。
\item 正确 6436/8000 = 0.8045，筛出 1564 条错误样本。
\item 只训 hard cases 后验证降至 \metric{0.834}{0.8346}{0.8254}。
\item 后续改为 hard cases 与常规 Task1/Task3 样本混合。
\end{itemize}
\end{block}
\end{column}
\end{columns}
\vspace{0.2em}
\scriptsize
\begin{verbatim}
Thought: <brief reasoning process>
Answer: <short answer>
\end{verbatim}
\end{frame}

\begin{frame}{最终方案：bbox prompt + hard case 混合训练}
\small
\begin{itemize}
\item 选择最稳定路径：统一使用 \texttt{bbox\_prompt}，D\_full，全参数长训练。
\item 先用原始模型在 8000 条训练数据上 bbox 推理，保留错误样本作为 hard cases。
\item 直接只训 hard cases 不稳定，因此采用全量训练与 hard case 混合。
\end{itemize}
\smalltab
\begin{tabular}{p{0.42\textwidth}cccp{0.21\textwidth}}
\toprule
训练设置 & Total & BBox & No-BBox & 观察 \\
\midrule
baseline：bbox\_prompt & 0.799 & 0.7972 & 0.8254 & 原始模型已显著提升 \\
8000 samples, bs16 lr1e-5 & 0.890 & 0.8922 & 0.8571 & 长训练提升 bbox subset \\
task1/3 hard 3128, bs16 lr5e-6 & 0.894 & 0.8954 & 0.8730 & hard case 混合有效 \\
shuffle 8000, bs64 lr5e-6 / best & \textbf{0.906} & \textbf{0.9082} & \textbf{0.8730} & 最优 checkpoint step 100 \\
\bottomrule
\end{tabular}
\vspace{0.35em}
\begin{block}{Leaderboard}
最终提交 Leaderboard 得分 \textbf{85.3}，排名第一。
\end{block}
\end{frame}

\begin{frame}{最终训练曲线与收益来源}
\small
\begin{columns}[T,onlytextwidth]
\begin{column}{0.56\textwidth}
\includegraphics[width=\textwidth,height=0.58\textheight,keepaspectratio]{\imgdir/task3_wandb.png}
\end{column}
\begin{column}{0.41\textwidth}
\begin{block}{关键结论}
\begin{itemize}
\item Task 3 的主要提升来自显式 grounding，而不是更复杂的 instruction 数据。
\item bbox prompt 把定位问题转化为模型可读的结构化上下文。
\item 标签/对象文本比单纯画框更稳定；坐标、颜色锚点提供辅助。
\item COT 对当前短答案精确匹配任务收益不大。
\end{itemize}
\end{block}
\end{column}
\end{columns}
\end{frame}

\section{总结}

\begin{frame}{总体结论}
\small
\begin{columns}[T,onlytextwidth]
\begin{column}{0.48\textwidth}
\begin{block}{回答实验问题}
\begin{itemize}
\item 视觉编码器并非总是最值得优先微调；当前规模下语言侧适配更有效。
\item Connector 负责跨模态对齐与输入格式适配，参数少但性价比高。
\item 全参数微调性能最高，但收益与参数量不线性，需要结合显存和时间权衡。
\end{itemize}
\end{block}
\end{column}
\begin{column}{0.48\textwidth}
\begin{block}{核心发现}
\begin{itemize}
\item Prompt 格式约束能显著改变 VQA 精确匹配结果。
\item Task2 合成 instruction 数据可行，但不天然迁移到 GQA 短答案分布。
\item 显式 bbox grounding 是最直接、最稳定的增益来源。
\item 最终 \texttt{bbox\_prompt + D\_full + hard case 混合} 达到 0.906 验证准确率。
\end{itemize}
\end{block}
\end{column}
\end{columns}
\end{frame}

\backmatter
\end{document}
